% ---------------------------------------------------------------------------
%  中文报告：复盘、周报、技术方案、实验报告都能用
%  数字卡片、时间线、提示框、代码块、决策表，换掉内容就是一份成品。
%  文中的日期与人名为示例；关于素帛的数字与做法取自它真实的设计。
% ---------------------------------------------------------------------------
\documentclass[UTF8,zihao=-4]{ctexart}

\usepackage[a4paper,textwidth=38\ccwd,top=2.6cm,bottom=2.6cm,headheight=16pt]{geometry}
\setmainfont{texgyrepagella}[
  Extension=.otf, UprightFont=*-regular, ItalicFont=*-italic,
  BoldFont=*-bold, BoldItalicFont=*-bolditalic]
\setmonofont{SourceCodePro}[Extension=.otf, UprightFont=*-Regular, BoldFont=*-Semibold, Scale=MatchLowercase]
\usepackage{xcolor}
\usepackage{booktabs}
\usepackage{tabularx}
\usepackage{enumitem}
\usepackage{fancyhdr}
\usepackage{listings}
\usepackage{tikz}
\usepackage[most]{tcolorbox}
\usepackage{hyperref}

% 配色：墨色做结构，绿色是唯一的强调，琥珀色只用来提示风险
\definecolor{ink}{HTML}{1F3A33}
\definecolor{accent}{HTML}{2F7D63}
\definecolor{amber}{HTML}{E0A63A}
\definecolor{mist}{HTML}{75837E}
\definecolor{paper}{HTML}{F4F7F5}
\hypersetup{colorlinks, allcolors=accent}

\linespread{1.4}
\setlist{itemsep=1pt, topsep=3pt, leftmargin=2em}
\ctexset{
  section    = {format=\Large\bfseries\color{ink}, name={,}, number=\arabic{section},
                aftername=\hspace{.8em}, beforeskip=3ex, afterskip=1.5ex},
  subsection = {format=\large\bfseries, beforeskip=2ex, afterskip=1ex},
}

% 斜杠加一个点：落款与页眉上的小标记。深底上用 \subomark*（反白）
\definecolor{snow}{HTML}{F7F8F7}
\definecolor{leaf}{HTML}{6FCF9F}
\NewDocumentCommand{\subomark}{s O{1em}}{{%
  \IfBooleanTF{#1}{\colorlet{slash}{snow}\colorlet{dot}{leaf}}{\colorlet{slash}{ink}\colorlet{dot}{accent}}%
  \tikz[x=#2/120, y=#2/120, baseline=-.1*#2]{%
    \draw[slash, line width=12*#2/120, line cap=round] (36,90) -- (62,32);
    \fill[dot] (82,34) circle (8);}}}

\pagestyle{fancy}
\fancyhf{}
\renewcommand{\headrulewidth}{0pt}
\fancyhead[L]{\small\color{mist}\subomark[.9em]\enspace 素帛 · 季度复盘}
\fancyhead[R]{\small\color{mist}内部资料}
\fancyfoot[C]{\small\color{mist}\thepage}

% 数字卡片：\kpi{数字}{说明}{变化}
\newcommand{\kpi}[3]{%
  \begin{tcolorbox}[enhanced, width=\linewidth, colback=paper, colframe=paper, arc=3pt,
    left=10pt, right=10pt, top=8pt, bottom=8pt, equal height group=kpi]
    {\fontsize{24}{28}\selectfont\bfseries\color{ink}#1}\par\vspace{2pt}
    {\small #2}\par
    {\footnotesize\color{accent}#3}
  \end{tcolorbox}}

% 提示框
\newtcolorbox{note}[1][要点]{enhanced, colback=white, colframe=accent, boxrule=0pt,
  leftrule=2.5pt, arc=0pt, left=12pt, top=6pt, bottom=6pt,
  fonttitle=\bfseries\color{accent}, coltitle=accent, colbacktitle=white,
  title=#1, attach title to upper={\quad}}
\newtcolorbox{risk}[1][风险]{enhanced, colback=amber!6, colframe=amber, boxrule=0pt,
  leftrule=2.5pt, arc=0pt, left=12pt, top=6pt, bottom=6pt,
  fonttitle=\bfseries, coltitle=amber!70!black, colbacktitle=amber!6,
  title=#1, attach title to upper={\quad}}

\lstset{
  basicstyle=\ttfamily\small, columns=fullflexible, breaklines=true,
  frame=none, backgroundcolor=\color{paper}, xleftmargin=10pt, framexleftmargin=10pt,
  aboveskip=8pt, belowskip=8pt,
  keywordstyle=\color{accent}\bfseries, commentstyle=\color{mist},
  stringstyle=\color{ink}, showstringspaces=false,
}

\begin{document}

% ---------------------------------------------------------------- 题头
\thispagestyle{empty}
\begin{tikzpicture}[remember picture, overlay]
  \fill[ink] (current page.north west) rectangle ([yshift=-3.6cm]current page.north east);
  \node[anchor=south west, inner sep=0pt] at ([xshift=2.6cm, yshift=-2.9cm]current page.north west)
    {\subomark*[3em]};
\end{tikzpicture}
\vspace*{1.6cm}

\noindent{\small\color{accent}\bfseries 复盘\quad\textperiodcentered\quad 2026 年第三季度}\par
\vspace{6pt}
\noindent{\fontsize{26}{32}\selectfont\bfseries\color{ink} 素帛，让灵感更自由\par}
\vspace{4pt}
\noindent{\large\color{mist} 一个离线 \LaTeX{} 工作台的第一个季度\par}
\vspace{14pt}
\noindent{\small
  \begin{tabularx}{\linewidth}{@{}l X l l@{}}
    {\color{mist}撰写} & 陈默（产品）、许嘉（工程） & {\color{mist}日期} & 2026 年 10 月 19 日 \\
  \end{tabularx}}
\vspace{10pt}

\begin{note}[一句话]
  我们把「装一个 \TeX{} 发行版」从几个 G、一个下午，缩短到一个两百来兆的安装包；
  用不到的宏包和字体一个也不带，缺什么，打开文档时再补。
\end{note}

% ---------------------------------------------------------------- 数字
\section{结果}

\noindent
\begin{minipage}[t]{.315\linewidth}\kpi{220\,MB}{安装包大小}{完整发行版约 8\,GB}\end{minipage}\hfill
\begin{minipage}[t]{.315\linewidth}\kpi{0}{编译所需的网络请求}{文档与编译都在本机}\end{minipage}\hfill
\begin{minipage}[t]{.315\linewidth}\kpi{1}{补齐缺包要点几下}{缺字清单，一键安装}\end{minipage}

\vspace{6pt}
三项目标全部达成。最让我们意外的是：只内置期刊、学位论文、简历模板里最常用的那一批宏包，
就能让大多数真实模板开箱即编；剩下的，交给「缺字清单」在第一次编译时一次补齐。

% ---------------------------------------------------------------- 背景
\section{我们要解决什么}

写 \LaTeX{} 的人通常面对三种不顺手：
\begin{itemize}
  \item \textbf{太重}：本地发行版动辄几个 G，装一次要一个下午；
  \item \textbf{太远}：云端编辑器依赖网络，网速慢时编译要排队，导出还可能收费；
  \item \textbf{太散}：编辑器、插件、发行版各管一段，出了错不知道该找谁。
\end{itemize}
「素帛」的回答是：文档和编译都留在本机，工具只带最小够用的那一部分。

% ---------------------------------------------------------------- 时间线
\section{过程}

\begin{center}
\begin{tikzpicture}[x=1cm, y=1cm, every node/.style={font=\small}]
  \draw[mist!60, line width=1pt] (0,0) -- (13.6,0);
  \foreach \x/\date/\what/\c in {
      0/7 月 8 日/立项/ink,
      3.4/8 月 2 日/编辑与预览跑通/ink,
      6.8/8 月 30 日/精简发行版/ink,
      10.2/9 月 21 日/缺字清单/ink,
      13.6/10 月 12 日/首个安装包/accent} {
    \fill[\c] (\x,0) circle (4pt);
    \node[above=6pt, color=mist] at (\x,0) {\date};
    \node[below=6pt, color=ink, font=\small\bfseries] at (\x,0) {\what};
  }
\end{tikzpicture}
\end{center}

\subsection{做对了什么}
\begin{enumerate}
  \item \textbf{先砍，再加。}从完整发行版出发，统计主流期刊、学位论文与简历模板
        实际用到的宏包，只保留出现频率最高的那一批。
  \item \textbf{把「缺」变成一个清单。}编译失败时不让用户读日志，而是直接列出缺的宏包，
        一键补齐：
\begin{lstlisting}[language=bash]
# 用户看到的是一张清单，背后只是一条命令
tlmgr install biblatex-gb7714-2015   # 缺 gb7714-2015.bbx
\end{lstlisting}
  \item \textbf{预置模板只带源码。}模板里引用的字体、宏包不随安装包分发，
        打开即编译，缺什么补什么。
\end{enumerate}

\subsection{做得不够的地方}
\begin{risk}[待改进]
  首次补包仍需联网；在网络受限的环境里，需要提供一个可离线导入的宏包合集。
\end{risk}

% ---------------------------------------------------------------- 决策
\section{关键决策}

\begin{center}
\small
\renewcommand{\arraystretch}{1.35}
\begin{tabularx}{\linewidth}{@{}l X X@{}}
  \toprule
  \textbf{问题} & \textbf{选择} & \textbf{理由} \\
  \midrule
  内置哪些宏包 & 按真实模板的使用频率选 & 体积与覆盖率的最好折中 \\
  中文字体     & 用系统自带字体           & 一款中文字体就有十几 MB \\
  编译引擎     & 只带 Xe\LaTeX{}          & 中文支持最好，维护一套就够 \\
  \bottomrule
\end{tabularx}
\end{center}

% ---------------------------------------------------------------- 下一步
\section{下一步}
\begin{enumerate}
  \item 离线宏包合集，覆盖无网络环境（许嘉，11 月底）；
  \item 模板分类与预置模板，让新用户五分钟内排出第一份文档（陈默，11 月中）；
  \item 编译错误的中文解释，覆盖最常见的 20 种报错（许嘉，12 月）。
\end{enumerate}

\vspace{.5em}
\hfill\subomark[1.4em]

\end{document}
